% ==============================================================================
% MÓDULO PRE-CAPÍTULOS: RESUMEN EN ESPAÑOL
% Archivo: pre_caps/05_resumen.tex
% Enfoque: 100% Calidad de Atención (Donabedian y Watson)
% ==============================================================================
\section*{RESUMEN}
\addcontentsline{toc}{section}{RESUMEN}

\noindent El presente proyecto de investigación tiene como objetivo general comprender el significado, estructura y esencia de la experiencia vivida por los usuarios respecto a la calidad de atención recibida en los consultorios del Centro de Salud Belén (Categoría I-3), provincia de Huamanga, Ayacucho, durante el año 2026. La investigación se fundamenta epistemológicamente en el paradigma interpretativo-naturalista bajo un diseño cualitativo fenomenológico (Edmund Husserl, Max van Manen y Roberto Hernández-Sampieri), articulado teóricamente con el enfoque del Proceso Interpersonal de Avedis Donabedian y la Teoría del Cuidado Humano de Jean Watson. El estudio explora la vivencia de la calidad del cuidado en sus cuatro existenciales del mundo de la vida (\textit{Lebenswelt}): relacionalidad (calidez del encuentro interpersonal, empatía de enfermería, acogida y diálogo en quechua chanka), temporalidad (oportunidad, ritmo pausado y tiempo dedicado a la consulta clínica), espacialidad (privacidad visual y auditiva, confort y confidencialidad en los consultorios) y corporalidad (respeto al pudor físico del cuerpo y contención humanizada del dolor y sufrimiento somático). La población de referencia comprende a los usuarios externos atendidos en el establecimiento asistencial (18,450 habitantes adscritos, con 68.2\% en situación de pobreza SISFOH y 72\% de bilingüismo activo quechua chanka-castellano). La selección de los participantes se realizará mediante muestreo cualitativo no probabilístico intencional, estableciendo una muestra estimada de 12 a 18 informantes clave regulada por saturación teórica de categorías. Las técnicas de recolección de vivencias incluirán la entrevista fenomenológica en profundidad apoyada en una guía semiestructurada bilingüe y la observación reflexiva registrada en notas de campo. El procesamiento testimonial se efectuará mediante el método de Colaizzi asistido por ATLAS.ti v9, transitando hacia la delimitación de categorías esenciales y divergentes de la calidad del cuidado. El rigor científico se garantizará conforme a los criterios de credibilidad, transferibilidad, consistencia y confirmabilidad (Guba y Lincoln, 1985), salvaguardando la reducción fenomenológica (\textit{epojé}) y los principios éticos de la Declaración de Helsinki.

\vspace{0.8cm}
\noindent \textbf{Palabras clave (DeCS):} Calidad de la Atención de Salud; Cuidado de Enfermería; Relaciones Interpersonales; Empatía; Humanización de la Atención; Investigación Cualitativa.

\newpage
